\section{An exact decomposition refinement for all Spin characters}
\label{sec:spin-parameter-optimization}

\paragraph{Status and separation from the withdrawn candidate.}
These are new statements D-OL-27--29. The complete author and root
whole-proof reviews and the separate /root/sixfolds whole-proof,
primary-source and exact-certificate audit have passed.  The withdrawn D-OL-25--26 is not a premise:
in particular the second line of LST (4.3.2) has exponent
$\sqrt N/6$, not $\sqrt{N/6}$, and is not used here.
The argument below instead uses the published decomposition and
restriction lemmas and proves a new constituent-counting inequality.
The scalar bound used throughout is the uniform $\gamma=19/20$;
no subsequent numerical assertion $\gamma_q\le1/7$ is needed.

All fields in this section have odd characteristic.  Write
\[
 \gamma=\frac{19}{20},\qquad
 \rho_0=\frac{99997}{150000},\qquad
 \beta_0=\frac16-\frac1{10^7},\qquad
 \delta_0=\frac16+\frac1{10^5},\qquad \rho_0=1-2\delta_0.
 \tag{SPO.1}
\]

\subsection*{D-OL-27. Counting trivial factors with their actual weights}

\begin{lemma}[D-OL-27]
Let $H=\langle K,h\rangle$ be finite, where
$K=K_1*\cdots*K_m$ is a normal central product and $h$ normalizes
each $K_i$.  Suppose the finite extensions $H_i$ and local operators
$h_i$ satisfy the action-compatibility hypotheses of LST Lemma~4.3.2,
and every $H_i$-character irreducible and nontrivial over $K_i$ has
absolute character ratio at $h_i$ at most $\gamma$, where $0<\gamma<1$.
For every complex $H$-module $A$ of dimension $a$,
\[
 \frac{|\operatorname{tr}(h|A)|}{a}
 \le\gamma^m+
       \frac{1-\gamma^m}{m}
          \sum_{i=1}^m\frac{\dim A^{K_i}}a.
 \tag{SPO.2}
\]
No irreducibility of $A|K$, independent realization of the centers,
or extension of every constituent to the full product is assumed.
\end{lemma}

\begin{proof}
Representations of a finite central product are the representations
of the direct product on which the central kernel acts trivially.
Its irreducibles therefore have tensor factors with the required
central compatibility.  For each subset $J\subseteq\{1,\ldots,m\}$,
let $A_J$ be the sum of the $K$-isotypic components whose trivial
factor set is exactly $J$.  Since $h$ normalizes every $K_i$,
it preserves whether that factor is trivial; hence $A_J$ is
an actual $H$-submodule.
LST Lemma~4.3.2, including its complete proof for arbitrary characters,
applies to the nontrivial factors outside $J$ and gives
\[
 |\operatorname{tr}(h|A_J)|
               \le\gamma^{m-|J|}\dim A_J.
 \tag{SPO.3}
\]
This remains valid when $h$ permutes irreducible $K$-constituents:
the published proof then gives zero trace on a nontrivial orbit.
When a constituent is invariant, cyclic extension and Schur's lemma
give the tensor operator up to a root of unity.  Thus neither
multiplicity spaces nor scalar phases are suppressed.
The exact primary source is
\cite[Lemma~4.3.2 and its full proof, pp.~1912--1913]{HOS:LST2011}.

The function $x\mapsto\gamma^{m-x}$ is convex.  Its chord between
$0$ and $m$ therefore bounds its value above:
\[
 \gamma^{m-k}\le\gamma^m+\frac{k}{m}(1-\gamma^m)
       \quad(0\le k\le m).
 \tag{SPO.4}
\]
For example its second derivative is
$(\log\gamma)^2\gamma^{m-x}>0$, which proves the stated chord
inequality.  On summing (SPO.3)--(SPO.4), use the exact identity
\[
 \sum_J |J|\dim A_J=\sum_{i=1}^m\dim A^{K_i}.
 \tag{SPO.5}
\]
This follows factor by factor: a nontrivial irreducible factor has
zero fixed space, while a trivial one fixes its whole component.
Equations (SPO.2)--(SPO.5) prove the lemma.
\end{proof}

\paragraph{A central-extension character bound used below.}
We record the precise cover interface rather than treating a genuine
Spin character as an ordinary orthogonal character.  If $Z$ is central
in a finite group $B$, $J=B/Z$, and $\eta\in\operatorname{Irr}(B)$,
then for every $b\in B$,
\[
                   |\eta(b)|^2\le |C_J(bZ)|.
 \tag{SPO.6}
\]
To prove this, let $\psi$ be the central character of $\eta$.
The central-character Fourier projector and ordinary column
orthogonality give
\[
 \sum_{\xi\in\operatorname{Irr}(B),\ \xi|_Z=\xi(1)\psi}
                 |\xi(b)|^2
  =\frac{|C_B(b)|}{|Z|}
        \sum_{z\in Z_b}\overline{\psi(z)},
 \quad Z_b=\{z\in Z:zb\text{ is conjugate to }b\}.
 \tag{SPO.7}
\]
The set $Z_b$ is a subgroup: it is the image of the inverse
image of $C_J(bZ)$ under the homomorphism recording the central
commutator with $b$.  Thus the last character sum is either
zero or $|Z_b|$.  Moreover
$|C_B(b)||Z_b|/|Z|=|C_J(bZ)|$.  The nonnegative summand
$|\eta(b)|^2$ in (SPO.7) proves (SPO.6).
This is a finite-group statement and has no faithfulness premise.

\subsection*{D-OL-28. A direct trace bound at restricted support 93}

\begin{theorem}[D-OL-28]
Let $q$ be odd, let $S=\operatorname{Spin}(V)$, and let
$\chi$ be any nontrivial irreducible complex character of $S$.
If $s\in S$ has natural restricted support
\[
 \operatorname{supp}_{\mu_2}(s)
      :=\min_{\varepsilon=\pm1}
             \operatorname{codim}\ker(\varphi(s)-\varepsilon I)
      \ge93,
 \tag{SPO.8}
\]
then $|\chi(s)|/\chi(1)<\rho_0$ and consequently
$\nu(\rho(s))>\delta_0\chi(1)$.
The representation may have any central kernel.
\end{theorem}

\begin{proof}
We state the exact published inputs and their domains before
choosing the new parameters.
For an embedded restriction subgroup $\operatorname{Spin}(U)$
with $\dim U=b\ge7$ and nondegenerate orthogonal complement,
the odd-characteristic orthogonal case in the complete proof
of LST Proposition~4.2.3 gives
\[
 \frac{\dim A^{\operatorname{Spin}(U)}}{\dim A}
       <\alpha_b(q):=\frac{q^{15/4-b}}{q-1}
 \tag{SPO.9}
\]
for every nontrivial irreducible complex module $A$ of the
ambient Spin group.  The precise source is
\cite[Proposition~4.2.3, pp.~1909--1912, its odd-orthogonal
case in Step~1 and the $s=1$ summation in Step~6, p.~1912]{HOS:LST2011}.
Step~6 explicitly bounds the sum by
$\sum_{n=b+1}^{\infty}q^{15/4-n}=q^{15/4-b}/(q-1)$.
This uses the published bound in that proof, not the weaker
all-classical bound $q^{8/3-b/2}$ in the proposition's statement.
For $b=7$ we use the Step~1 telescoping argument itself:
every term has index $n\ge b+1=8$, within the proved range
$n\ge7$ of the odd-orthogonal bound.  Thus this case does not
extend the proposition's displayed all-classical $N\ge8$
statement without proof.  Both ambient parities are allowed.
It is a statement about
simply connected groups, so is not restricted to ordinary
$\Omega$ characters.

The decomposition and extension inputs are LST
Proposition~4.3.3 and the complete Steps~1--3 of Theorem~4.3.6,
printed pp.~1914--1916 and pp.~1918--1922.
An admissible decomposition is an orthogonal sum of nondegenerate
$s$-stable natural subspaces with determinant-one restrictions.
The resulting subgroups $K_i=\operatorname{Spin}(U_i)$ form a
central product.  Step~1 constructs finite extensions inside
the special Clifford groups and local operators $s_i$ having
exactly the same conjugation on $K_i$ as $s$.
For $\dim U_i\ge5$, their derived groups are quasisimple and
simply connected, and the Gluck bound (4.1.1) applies to
characters irreducible and nontrivial over $K_i$.
In particular its ratio is at most $\gamma=19/20$ for every
odd $q$.  These are the exact Spin-cover and admissibility
interfaces stated and proved in Step~1, pp.~1918--1919.
Neither a natural outer-automorphism classification at dimension
eight nor an assumption excluding triality is used: only the
actual Clifford block actions constructed there are needed.

We make explicit the large-indecomposable-block estimate, including
the dimension-seven endpoint.  For an admissibly indecomposable
natural block $U$ of dimension $b\ge8$, the final paragraph of the
odd-orthogonal proof of Proposition~4.3.3 gives
\[
 |C_{O(U)}(\bar s_U)|<2q^b,\qquad
 d(\operatorname{Spin}(U))>\tfrac89q^{b-3}.
 \tag{SPO.10}
\]
Both inequalities are stated there for $b\ge7$, not just $b\ge12$.
The definition of $d$ as the minimal nontrivial complex
representation degree is p.~1908; the precise inequalities are
on p.~1916, with $L(U)=\operatorname{Spin}(U)$ defined on p.~1913.
Let $B=\langle K_U,s_U\rangle$ be the special-Clifford extension.
Its natural-map kernel is central and its image lies in $SO(U)$.
For every $\eta\in\operatorname{Irr}(B)$ irreducible and nontrivial
over $K_U$, the central-quotient bound (SPO.6) therefore gives
\[
 \frac{|\eta(s_U)|}{\eta(1)}
 <\frac{9\sqrt2}{8}q^{3-b/2}
 \le\frac{3\sqrt2}{8}<\frac9{16}\quad(b\ge8).
 \tag{SPO.11}
\]
The last strict inequality is $\sqrt2<3/2$.
This treats the genuine Spin characters without identifying them
with ordinary orthogonal characters.

Here is the additional endpoint proof when $b=7$.
Write the natural element as a commuting semisimple-unipotent
product.  The complete odd-orthogonal classification in
Proposition~4.3.3, pp.~1915--1916, shows that a non-$\pm1$
primary summand has even dimension and gives an admissible
orthogonal summand.  It cannot occupy all seven dimensions.
The minus-one primary space has even dimension because the
natural element has determinant one; both the plus and minus
restrictions have determinant one.  Admissible indecomposability
therefore forces the minus-one primary space to be zero.
The element is unipotent.  The orthogonal unipotent decomposition
in that same proof consists of blocks $J_{2k+1}$ or $2J_{2k}$,
each with determinant one.  In dimension seven, admissible
indecomposability thus forces a single regular block $J_7$.

We compute its full orthogonal centralizer directly.
The general-linear commutant is
$R=\mathbf F_q[x]/(x^7)$, with $T=1+x$ representing $J_7$.
The adjoint involution is $\sigma(T)=T^{-1}$, hence
$\sigma(x)=-x+x^2-\cdots+x^6$.
Its diagonal on the augmentation ideal $(x)$ in the basis
$x,x^2,\ldots,x^6$ is $(-1)^j$.
Since $q$ is odd, the involution is semisimple and its fixed
space on that six-dimensional ideal has dimension three.
Write $U_R=1+(x)$ for the commutative odd-order unit group,
distinguishing it from the natural block $U$.  Consequently,
\[
 |U_R|=q^6,\qquad |U_R^\sigma|=q^3,\qquad
 |\ker(U_R\longrightarrow U_R^\sigma,\ u\mapsto u\sigma(u))|=q^3.
 \tag{SPO.11a}
\]
Indeed the norm is onto: squaring is a bijection on $U_R$, and
the unique square root of any $\sigma$-fixed unit is itself
$\sigma$-fixed.  Multiplying by the constant units shows that
$|C_{O(U)}(J_7)|=2q^3$, since a constant has norm one exactly
when it is $1$ or $-1$.
The same minimum-degree bound from (SPO.10), now at $b=7$,
and (SPO.6) imply
\[
 \frac{|\eta(s_U)|}{\eta(1)}
 <\frac{9\sqrt2}{8}q^{-5/2}
 \le\frac{\sqrt{2/3}}8<\frac18.
 \tag{SPO.11b}
\]
No assertion about arbitrary outer automorphisms of dimension-seven
or dimension-eight groups is needed.

Refine the natural space into admissibly indecomposable blocks.
If some block has dimension at least seven, its fixed fraction
is bounded by (SPO.9).  The complement of the fixed space is
stable under $s$; the one-factor case of LST Lemma~4.3.2 applies
to its actual local extension.  At $b=7$,
$\alpha_7(3)<1/54$; at $b\ge8$, $\alpha_b(q)\le\alpha_8(3)<1/162$.
Thus in the two cases the full trace ratio is respectively less than
\[
 \frac18+\frac1{54}=\frac{31}{216}<\frac35,
 \qquad
 \frac9{16}+\frac1{162}=\frac{737}{1296}<\frac35<\rho_0.
 \tag{SPO.12}
\]
If the block is the entire ambient group its fixed fraction is zero.
This is Case A, with $N_1=7$.

It remains to handle the case where every indecomposable block
has dimension at most $6$.
Call a block good when its local operator is noncentral in
the corresponding extension, and bad otherwise.
The precise classification and pairing from the published
Step~3, pp.~1921--1922, applies as follows.
A bad block has dimension at most two and is either scalar
$1$ or $-1$, or is a two-dimensional rotation with eigenvalue
pair $\{\lambda,\lambda^{-1}\}$, $\lambda\ne\pm1$.
Pair different bad types, and then pair equal remaining
rotation types.  Every such pair becomes good and has
dimension at most four.  The unpaired remainder is either
at most one two-dimensional rotation block or consists
of scalar blocks of a single sign.
Consequently the total dimension $\Sigma_g$ of the good
blocks satisfies
\[
              \Sigma_g\ge\operatorname{supp}_{\mu_2}(s)-2.
 \tag{SPO.13}
\]
For a rotation remainder this follows by subtracting at most
two from the full dimension.  For a single scalar sign,
the bad subspace lies in that sign's eigenspace, so the
stronger bound without minus two holds.  This explains
directly why the restricted scalar set $\mu_2$ suffices
here, without importing an incorrectly transcribed exponent.
Every good block now has dimension in $\{3,4,5,6\}$:
in dimensions one and two the determinant-one isometry group
and its special-Clifford extension are abelian, so such a block
is bad.  The above pairings create only dimensions three or four.

We use a sharper elementary packing argument, with $m_0=N_2=8$.
Any multiset of integers in $\{3,4,5,6\}$ of total at least $91$
contains eight disjoint groups each of sum at least eight.
To prove this, choose a packing with the maximum number of such
groups, then make each group inclusion-minimal by returning
removable items to the remainder.  Each minimal group has sum
at most twelve; a sum-twelve group must be exactly $6+6$.
The remainder has sum at most seven.  If its sum is seven,
it is $3+4$.  In that case no $6+6$ group can occur, since
replacing it and the remainder by $6+3$ and $6+4$ would increase
the number of groups.  Therefore a packing with at most seven
groups has total at most
\[
 \max\{7\cdot12+6,\ 7\cdot11+7\}=90.
 \tag{SPO.14}
\]
This proves the assertion.  The threshold is sharp for this
packing problem, as fifteen items of size six have total ninety
and yield at most seven groups.

By (SPO.8)--(SPO.13), the good blocks have total at least
$93-2=91$.  Apply this packing, and append all unused blocks,
including the bad blocks, to the eighth group.  The resulting
admissible decomposition has eight subspaces of dimension at
least eight, each containing a good block and hence carrying a
noncentral local operator.  This uses the published classification
and admissibility interfaces, but the packing improvement is
proved here.  No real-valued parameter or square-root exponent
from the numerical conclusion of LST is chosen.

Apply D-OL-27 and (SPO.9) to these eight actual Spin
subgroups.  Monotonicity in their dimensions and in $q$
gives the same upper bound for each fixed fraction,
\[
 \alpha_8(q)\le\alpha_*:=\frac{3^{-17/4}}2,
 \qquad
 \frac{|\chi(s)|}{\chi(1)}
       <c+(1-c)\alpha_*,\quad c=(19/20)^8.
 \tag{SPO.15}
\]
All constants are certified by a short positive-integer
comparison.  Direct rational subtraction gives
\[
 u:=\frac{\rho_0-c}{1-c}
       =\frac{247774877}{25849310877}>0,
 \qquad
 16\cdot3^{17}\cdot247774877^4>25849310877^4.
 \tag{SPO.16}
\]
The second inequality is exactly $u^4>\alpha_*^4$.
Since both quantities are positive, $u>\alpha_*$;
(SPO.15)--(SPO.16) give the required ratio $<\rho_0$.
The finite certificate checks this integer inequality
directly.  Finally $|\chi(s)|\ge\chi(1)-2\nu(\rho(s))$
after rotating a largest-eigenspace eigenvalue, so the
support exceeds $\delta_0\chi(1)$.
This proves the theorem for every central kernel.
\end{proof}

\subsection*{D-OL-29. The resulting actual-source parameter bound}

\begin{theorem}[D-OL-29]
An actual rank-nine six-dimensional multiplicity source cannot
have a nonempty irreducibly acting normal central product of
represented odd-characteristic Spin factors for which every
odd natural dimension is at least $187$ and every even natural
dimension is at least $188$.
Every nontrivial irreducible complex factor module and every
represented central kernel is allowed, together with all
compatible scalars, outer automorphisms and factor permutations.
No selfduality or transitivity is required.
In the transitive-core problem the remaining necessary natural
dimension is therefore at most $185$ in odd dimension and
at most $186$ in even dimension.
\end{theorem}

\begin{proof}
We identify the changed quantitative interfaces in the
already complete D-OL-19--24 normalizer proof.
The Clifford lift and actual represented-kernel preservation
need only natural dimension at least $22$.
All fixed-hyperbolic-pair, quadratic-field and arbitrary
semilinear descent cases give support at least $3a/16$ or
$7a/32$, except the rank-one reflection case.
Their radical complements have dimension at least $d-4$,
so the elementary lower-dimensional conditions remain valid.
An inner prime-order case without a fixed hyperbolic pair
has restricted natural support at least
$\lceil(d-2)/2\rceil$, which is at least $93$ at both
new thresholds.  D-OL-28 replaces the previous invocation
of the first line of LST(4.3.2) there.
In the odd-dimensional all-orders reduction, the square of
a possible $2^s$-order root with $s\ge3$ lies in $\Omega$
and has restricted support at least $d-1>93$;
D-OL-28 applies to its Spin lift as well.
The remaining fourth-order case is D-OL-22 and requires
only $d\ge25$.
The complete even-dimensional determinant reduction in
D-OL-20 is unchanged.  Thus no extra natural-dimension
condition has been omitted.

Here is the sharper rank-one estimate in full, without
using the withdrawn D-OL-25 as an input.
In the nondegenerate radical dual of dimension $b\ge43$,
a nonzero norm orbit has size at least
$q^{b-1}-q^{b/2}-1$ by the elementary quadratic Fourier
sum.  The fixed hyperplane of a reflection is nondegenerate
and has at most $q^{b-2}+q^{(b-1)/2}$ vectors of that norm.
Therefore its fixed fraction is at most
\[
 c_q^*:=\frac{q^{-1}+q^{-20}}{1-q^{-20}-q^{-40}}.
 \tag{SPO.17}
\]
Indeed, after division by $q^{b-1}$ the numerator error
is $q^{-(b-1)/2}\le q^{-21}$ and the denominator errors
are $q^{1-b/2}\le q^{-20}$ and $q^{1-b}\le q^{-40}$.
The underlying Gauss bound follows by adding the $q-1$
nontrivial additive-character terms, each of magnitude
$q^{b/2-1}$; it applies also to norm zero, with the
possible zero vector removed in the denominator.
The fraction in (SPO.17) decreases with $q$.
The exact rational endpoints are
\[
           c_3^*<1-4\beta_0,\qquad c_5^*<1/4.
 \tag{SPO.18}
\]
The lifted opposite-radical fixed space has dimension
at most $a/2$, and nonfixed involution weight orbits
have length two.  Thus the rank-one support is at least
$(1-c_q^*)a/4$, greater than $\beta_0a$ at $q=3$ and
greater than $3a/16>\delta_0a$ at $q\ge5$.

It follows, with every projective order retained as in
D-OL-20 and D-OL-23, that every nonscalar compatible
normalizer has support at least $\beta_0a$.
Any operator with support at most $\delta_0a$ has
$q=3$ and natural projective lift a reflection of
projective order two.  Actual squares may still be
central scalars, as explicitly permitted in those proofs.

For every nontrivial factor the lifted-radical weight
orbit degree bound gives, with $b=d-2\ge185$,
\[
 a\ge q^{b-1}-q^{\lceil b/2\rceil}-1
      \ge(q-1)q^{b-2}-1
      \ge2\cdot3^{183}-1>10^6.
 \tag{SPO.19}
\]
This uses only a nontrivial weight and its full inverse-image
Levi orbit, so includes genuine Spin modules.
Write $e$ for the product of the factor degrees; then $e>10^6$.

Call an actual inertia good when it moves at least
$\delta_0e$ dimensions.  Every genuine inertia moves
at least $\beta_0e$; nonidentity scalars move all $e$,
and nontrivial factor cycles have support at least $e/4$.
A bad inertia is therefore factor preserving, with
every active factor a projective reflection over $\mathbf F_3$.
Two such active factors have support at least
$2\beta_0(1-\beta_0)e>e/4$, by the exact two-eigenspace
convolution calculation.  Thus a bad inertia has exactly
one active factor, independent of its scalar phase.

The actual intermediate-cover Euler and moduli arguments
force a sphere base with at least twelve genuine marks,
total moving rank $2e+6$, and a generating product-one tuple.
The exact inequalities
\[
 13\beta_0-2>1/7,\quad
 11\beta_0+\delta_0-2=89/10^7,\quad
 6\cdot10^7/89<10^6
 \tag{SPO.20}
\]
exclude thirteen marks or even one good inertia among twelve.
Exactly twelve bad inertias remain.
Their projections are trivial on factors over $q\ne3$,
which is already incompatible with generation.
On a factor over $\mathbf F_3$, at most twelve natural
reflections represent their projective images.
The common perpendicular of their normals has dimension
at least $d-12>0$ and is fixed by every reflection word.
Every element of the natural $\Omega$ is projectively
such a word and therefore acts by a ground-field scalar
on that common subspace.  Perfectness kills this scalar
character, contradicting the absence of a nonzero natural
fixed space.  In even dimension this is argued through
$P\Omega$ exactly as in D-OL-21; no splitting of either
the natural or represented Spin center is introduced.
The contradiction proves the theorem.
\end{proof}

\paragraph{Exact evidential boundary.}
The weighted trivial-factor estimate, central-extension column
bound, regular-seven-block centralizer count and sharp small-block
packing are proved explicitly above.  The last two refinements
were contributed by the root collaborator and checked here; this
is a joint derivation, not an independent-origin audit.
The decomposition and subgroup estimates are invoked with
their complete published hypotheses and exact proof locators.
The finite certificate checks all constants and packing
endpoints, and does not replace those universal arguments.
This yields necessary remaining parameter bounds and
asserts no actual source in the surviving dimensions.
